\item En la figura P7.22, el cambio en energía interna de un gas que se lleva de $A$ a $C$ a lo largo de la trayectoria azul tiene el código QR) es $+800\text{ J}$. El trabajo efectuado sobre el gas a lo largo de la trayectoria roja $ABC$ es $-500\text{ J}$. (a) ¿Cuánta energía se debe agregar al sistema por calor a medida que va de $A$ a $B$ a $C$? (b) Si la presión en el punto $A$ es cinco veces la del punto $C$, ¿cuál es el trabajo efectuado sobre el sistema de $C$ a $D$? (c) ¿Cuál es la energía que se intercambia con los alrededores por calor a medida que el gas va de $C$ a $A$ a lo largo de la trayectoria verde? (d) Si el cambio en energía interna al ir del punto $D$ al punto $A$ es $+500\text{ J}$, ¿cuánta energía se debe agregar al sistema por calor a medida que va del punto $C$ al punto $D$?
